% Related Work — v0.1 skeleton.
%
% Target: 1-2 pages. Positions PrimitiveHub against three lines of
% prior work: software supply-chain security, prompt-injection
% defense, and AV-vendor methodology.

\section{Related Work}
\label{sec:related}

% Subsection 1 — software supply chain.
\subsection{Software supply-chain provenance}

% Claim 1.1 — SLSA + in-toto + Sigstore.
% TODO: SLSA \cite{SLSA:2023:provenance-spec} formalises build-
% provenance levels. in-toto \cite{in-toto:2016:supply-chain}
% provides the predicate-based attestation framework. Sigstore
% \cite{Sigstore:2021:keyless-signing} provides keyless signing via
% OIDC. PrimitiveHub's Trust Bundle inherits the predicate format
% from in-toto and the signing layer from Sigstore; the addition
% is the SEMANTIC validation layer (Top 10) that runs BEFORE the
% signing step.

% Claim 1.2 — incident catalog.
% TODO: cite the chain of supply-chain incidents the threat model
% draws from — SolarWinds \cite{SolarWinds:2020:supply-chain},
% Codecov \cite{Codecov:2021:bash-uploader}, ua-parser-js
% \cite{ua-parser-js:2021:npm-typosquat}, eslint-config-prettier
% \cite{eslint-config-prettier:2024:maintainer-compromise}, xz-utils
% CVE-2024-3094 \cite{xz-utils:2024:CVE-2024-3094}, and the
% Shai-Hulud npm worm \cite{Shai-Hulud:2025:npm-worm}. Each
% motivates one or more CCI categories.

% Subsection 2 — prompt-injection defense.
\subsection{Prompt-injection defense}

% Claim 2.1 — what the field has tried.
% TODO: prior work has focused on runtime defenses — input/output
% filtering, prompt rewriting, instruction-hierarchy training,
% RAG-retrieval sanitization. These operate at LLM execution time
% and don't prevent the attack from reaching the model.

% Claim 2.2 — what's different.
% TODO: PrimitiveHub's scanners operate at REGISTRY PUBLISH TIME.
% The artifact never reaches the consuming agent unless it has
% passed the Top 10. This is the same shift package managers made
% from "warn on install" to "block install of known-bad" — moving
% the defense earlier in the pipeline.

% Subsection 3 — OWASP LLM Top 10.
\subsection{OWASP Top 10 for LLM applications}

% Claim 3.1 — the relationship.
% TODO: OWASP \cite{OWASP:2025:llm-top10} defines a 10-category
% taxonomy for LLM application security. PrimitiveHub's Top 10
% overlaps with OWASP's LLM01 (Prompt Injection), LLM02 (Insecure
% Output Handling), LLM03 (Training Data Poisoning), and LLM05
% (Supply Chain Vulnerabilities). The crosswalk is documented in
% docs/owasp-llm-top10-crosswalk.md.

% Claim 3.2 — the scope boundary.
% TODO: PrimitiveHub covers the supply-chain + injection slice that
% can be detected statically on agent-context content. OWASP's
% runtime categories (LLM02 Insecure Output Handling, LLM04 Model
% Denial of Service, LLM06 Sensitive Information Disclosure
% post-inference) require different defense layers — typically
% provider-side or proxy-side, not registry-side.

% Subsection 4 — AV-vendor methodology.
\subsection{Antivirus industry methodology}

% Claim 4.1 — the parallel.
% TODO: the AV industry has spent two decades formalising
% scanner-as-product. The product-family map (signature DB / FP
% portal / allowlist DB / on-access guard / boot-time scan / HIPS /
% cloud lookup / provenance) is documented in
% \cite{AntivirusTaxonomy:2010s:industry-practice}. PrimitiveHub
% mirrors this map row-for-row (docs/decisions/0030 § 3).

% Claim 4.2 — what carries over.
% TODO: append-only FP/FN ledgers, four-path triage, quarterly
% re-review, per-detector versioning, corpus-as-ground-truth. These
% norms are well-evidenced in AV practice; PrimitiveHub adopts them
% for agent-context AV.

% Claim 4.3 — what's new.
% TODO: agent-context primitives are MARKDOWN with SEMANTIC content.
% AV detectors operate on binary executables. The detection
% techniques — regex pools with markdown-context awareness, UTS #39
% skeleton comparison for vendor names, codepoint-class density
% checks — are PrimitiveHub-specific even when the surrounding
% methodology is AV-industry-standard.

% Subsection 5 — Unicode security.
\subsection{Unicode-based attacks}

% Claim 5.1 — UTS #39 + confusables.
% TODO: cite UTS #39 \cite{UTS39:2024:unicode-confusables} and the
% confusables.txt distribution. The skeleton method we use for
% T7 vendor-homoglyph detection is the standard one; what's new is
% applying it to PRIMITIVE-vendor identifiers at publish time.

% Claim 5.2 — the npm typosquat tradition.
% TODO: confusables-based phishing detection
% \cite{Confusables:2017:phishing-detection} has a decade of
% practice in the browser-security space. PrimitiveHub brings it
% into the primitive-registry space.

% Subsection 6 — measurement methodology.
\subsection{Measurement methodology}

% Claim 6.1 — precision/recall/F1.
% TODO: standard binary-classification evaluation as documented in
% \cite{Powers:2011:evaluation-precision-recall}. Per-detector
% macro-averaging chosen because each CCI category is a separate
% defense (no class-imbalance averaging argument).

% Claim 6.2 — holdout.
% TODO: the train/test split discipline
% \cite{HoldoutValidation:1995:tradition} is industry-standard
% since the 1990s. PrimitiveHub's 80/20 ratio + per-category
% balance follows the convention.
